\chapter{Code and Reproducibility}
\label{app:code}

\section{Repository Structure}
\label{sec:app-repos}

% Structure of the two repositories with full directory paths
% A table or directory tree
% Brief explanation of why the project is split into two repositories
% (different development lifecycles, cleaner commit histories)

\section{Computational Environment}
\label{sec:app-env}

% Conda environment name: amoc-thesis
% Python version, Jupyter version
% List of core packages with fixed versions:
% numpy, xarray, netCDF4, scikit-learn, PyTorch, matplotlib, plotly
% Dedicated Jupyter kernel

\section{Data Access and Preprocessing}
\label{sec:app-data}

% Links to the RAPID and Copernicus datasets
% Explanation of why the datasets are not included in Git
% (file size, licensing constraints)
% Scripts/commands for downloading the datasets
% Expected directory: /03_Data/raw/

\section{Analytical Pipeline}
\label{sec:app-pipeline}

% Execution order of notebooks NB01–NB05
% Output produced by each notebook
% Expected execution time
% Instructions for reproducing the main figures

\section{Bibliographic Management}
\label{sec:app-biblio}

% Zotero + Better BibTeX
% Citation key convention: [auth:lower][year]
% Automatic export workflow to references.bib

\section{Thesis Compilation}
\label{sec:app-latex}

% TeX Live 2026, biblatex + biber
% VS Code + LaTeX Workshop
% Build command: latexmk -pdf main.tex
% .gitignore file for build artifacts

\section{Version Control Workflow}
\label{sec:app-git}

% Commit message conventions
% Protected main branch and feature branches
% How to restore a previous project state
% Backup strategy